\chapter*{参考文献}
\markboth{参考文献}{参考文献}
\addcontentsline{toc}{chapter}{参考文献}
\sourcepage{243--245}
以下保留原书文献编号、作者姓名及出版资料；中文题名后列原题名，便于检索。
\section*{(a) 单复变函数论}
\begin{description}[style=nextline,leftmargin=2em]
\item[{[1]}] Ahlfors, L. V.：《复分析》（Complex Analysis）。New York: McGraw-Hill, 1966。
\item[{[2]}] Behnke, H.; Sommer, F.：《单复变解析函数论》（Theorie der analytischen Funktionen einer komplexen Veränderlichen），第3版。Berlin--Heidelberg--New York: Springer, 1965。
\item[{[3]}] Cartan, H.：《单复变与多复变解析函数初等理论》（Elementary Theory of Analytic Functions of One or Several Complex Variables）。Reading, MA: Addison-Wesley, 1963。
\item[{[4]}] Hurwitz, A.; Courant, R.：《函数论》（Funktionentheorie），附录由 H. Röhrl 撰写，第4版。Berlin--Heidelberg--New York: Springer, 1964。
\item[{[5]}] Lang, S.：《复分析》（Complex Analysis）。Reading, MA: Addison-Wesley, 1977。
\end{description}
Behnke--Sommer 和 Hurwitz--Courant 两书中有相当大的篇幅讨论黎曼曲面。
\section*{(b) 黎曼曲面}
\begin{description}[style=nextline,leftmargin=2em]
\item[{[10]}] Ahlfors, L. V.; Sario, L.：《黎曼曲面》（Riemann Surfaces）。Princeton, NJ: University Press, 1960。
\item[{[11]}] Chevalley, C.：《单变量代数函数论导引》（Introduction to the Theory of Algebraic Functions of One Variable）。American Mathematical Society, 1951。
\item[{[12]}] Farkas, H. M.; Kra, I.：《黎曼曲面》（Riemann Surfaces）。New York--Heidelberg--Berlin: Springer, 1980。
\item[{[13]}] Guenot, J.; Narasimhan, R.：《黎曼曲面论导引》（Introduction à la Théorie des Surfaces de Riemann）。Monographies de l'Enseignement Mathématique 23, Genève, 1976。
\item[{[14]}] Gunning, R. C.：《黎曼曲面讲义》（Lectures on Riemann Surfaces）。Princeton Mathematical Notes 2 (1966)。
\item[{[15]}] Gunning, R. C.：《黎曼曲面上的向量丛讲义》（Lectures on Vector Bundles over Riemann Surfaces）。Princeton Mathematical Notes 6 (1967)。
\item[{[16]}] Gunning, R. C.：《黎曼曲面讲义：雅可比簇》（Lectures on Riemann Surfaces: Jacobi Varieties）。Princeton Mathematical Notes 12 (1972)。
\item[{[17]}] Nevanlinna, R.：《单值化》（Uniformisierung）。Berlin--Heidelberg--New York: Springer, 1953。
\item[{[18]}] Pfluger, A.：《黎曼曲面论》（Theorie der Riemannschen Flächen）。Berlin--Heidelberg--New York: Springer, 1957。
\item[{[19]}] Serre, J.-P.：《代数群与类域》（Groupes Algébriques et Corps de Classes）。Paris: Hermann, 1959。
\item[{[20]}] Springer, G.：《黎曼曲面导论》（Introduction to Riemann Surfaces）。Reading, MA: Addison-Wesley, 1957。
\item[{[21]}] Weyl, H.：《黎曼曲面的概念》（The Concept of a Riemann Surface）。Reading, MA: Addison-Wesley, 1955。
\end{description}
Weyl 的书于1913年首次以德文出版，是第一部以现代方式处理黎曼曲面的著作，至今仍值得一读；其中还有许多十九世纪文献的引述。Chevalley 和 Serre 的书从代数的观点研究黎曼曲面。
\section*{(c) 多复变分析与复流形}
\begin{description}[style=nextline,leftmargin=2em]
\item[{[30]}] Andreian Cazacu, C.：《多复变函数论》（Theorie der Funktionen mehrerer komplexer Veränderlichen）。Berlin: Deutscher Verlag der Wissenschaften, 1975。
\item[{[31]}] Grauert, H.; Remmert, R.：《施泰因空间理论》（Theory of Stein Spaces）。Berlin--Heidelberg--New York: Springer, 1979。
\item[{[32]}] Gunning, R. C.; Rossi, H.：《多复变解析函数》（Analytic Functions of Several Complex Variables）。Englewood Cliffs, NJ: Prentice-Hall, 1965。
\item[{[33]}] Hirzebruch, F.：《代数几何中的拓扑方法》（Topological Methods in Algebraic Geometry）。Berlin--Heidelberg--New York: Springer, 1966。
\item[{[34]}] Hörmander, L.：《多复变复分析导论》（An Introduction to Complex Analysis in Several Variables），第2版。Amsterdam: North-Holland, 1973。
\item[{[35]}] Wells, R. O.：《复流形上的微分分析》（Differential Analysis on Complex Manifolds）。Englewood Cliffs, NJ: Prentice-Hall, 1973；第2版，Berlin--Heidelberg--New York: Springer, 1980。
\end{description}
\section*{(d) 拓扑、可微流形与泛函分析}
\begin{description}[style=nextline,leftmargin=2em]
\item[{[40]}] Bröcker, T.; Jänich, K.：《微分拓扑导论》（Einführung in die Differentialtopologie）。Heidelberger Taschenbücher 143, Berlin--Heidelberg--New York: Springer, 1973。
\item[{[41]}] Godement, R.：《代数拓扑与层论》（Topologie Algébrique et Théorie des Faisceaux）。Paris: Hermann, 1958。
\item[{[42]}] Massey, W. S.：《代数拓扑导论》（Algebraic Topology: an Introduction）。New York--Heidelberg--Berlin: Springer, 1967。
\item[{[43]}] Narasimhan, R.：《实流形与复流形上的分析》（Analysis on Real and Complex Manifolds）。Amsterdam: North-Holland, 1968。
\item[{[44]}] Schaeffer, H.：《拓扑向量空间》（Topological Vector Spaces）。New York: MacMillan, 1966。
\item[{[45]}] Schubert, H.：《拓扑学》（Topologie）。Stuttgart: Teubner, 1964。
\item[{[46]}] Seifert, H.; Threlfall, W.：《拓扑学教程》（Lehrbuch der Topologie）。Leipzig: Teubner, 1934；Chelsea, 1947 重印；英译本，Academic Press, 1980。
\item[{[47]}] Trèves, F.：《拓扑向量空间、分布与核》（Topological Vector Spaces, Distributions and Kernels）。New York--London: Academic Press, 1967。
\item[{[48]}] Warner, F.：《可微流形与李群基础》（Foundations of Differentiable Manifolds and Lie Groups）。Glenview, IL--London: Scott, Foresman, 1971。
\end{description}
\section*{(e) 专题}
\begin{description}[style=nextline,leftmargin=2em]
\item[{[50]}] Ahlfors, L. V.：闭黎曼曲面空间的复解析结构（The Complex Analytic Structure of the Space of Closed Riemann Surfaces）。载 Analytic Functions, Princeton, NJ: University Press, 1960, 45--66。
\item[{[51]}] Behnke, H.; Stein, K.：黎曼曲面上解析函数的展开（Entwicklungen analytischer Funktionen auf Riemannschen Flächen）。Math. Ann. 120 (1948), 430--461。
\item[{[52]}] Bieberbach, L.：《以函数论为基础的常微分方程理论》（Theorie der gewöhnlichen Differentialgleichungen auf funktionentheoretischer Grundlage dargestellt），第2版。Berlin--Heidelberg--New York: Springer, 1965。
\item[{[53]}] Cartan, H.：复解析簇与上同调（Variétés analytiques complexes et cohomologie）。Colloque sur les fonctions de plusieurs variables, CBRM, Bruxelles, 1953, 41--55。
\item[{[54]}] Florack, Herta：非闭黎曼曲面上的正则函数和亚纯函数（Reguläre und meromorphe Funktionen auf nicht geschlossenen Riemannschen Flächen）。Schriftenreihe Math. Inst. Univ. Münster 1 (1948)。
\item[{[55]}] Malgrange, B.：卷积偏微分方程解的存在性与逼近（Existence et approximation des solutions des équations aux dérivées partielles à convolution）。Ann. Inst. Fourier 6 (1955/56), 271--355。
\item[{[56]}] Meis, T.：紧黎曼曲面具体实现的最小叶数（Die minimale Blätterzahl der Konkretisierungen einer kompakten Riemannschen Fläche）。Schriftenreihe Math. Inst. Univ. Münster 16 (1960)。
\item[{[57]}] Röhrl, H.：线性微分方程论中的黎曼--希尔伯特问题（Das Riemann-Hilbertsche Problem der Theorie der linearen Differentialgleichungen）。Math. Ann. 133 (1957), 1--25。
\item[{[58]}] Ross, S. L.：《微分方程》（Differential Equations）。New York: Blaisdell, 1964。
\item[{[59]}] Serre, J.-P.：一般理论在若干整体问题中的应用（Applications de la théorie générale à divers problèmes globaux）。Séminaire H. Cartan, E.N.S., Paris, 1951/52, Exposé 20；重印本，Benjamin, 1967。
\item[{[60]}] Serre, J.-P.：关于完全连续映射的两个定理（Deux théorèmes sur les applications complètement continues）。Séminaire H. Cartan, E.N.S., Paris, 1953/54, Exposé 16；重印本，Benjamin, 1967。
\item[{[61]}] Serre, J.-P.：与施泰因簇有关的若干整体问题（Quelques problèmes globaux relatifs aux variétés de Stein）。Colloque sur les fonctions de plusieurs variables, CBRM, Bruxelles, 1953, 57--68。
\item[{[62]}] Stein, K.：具有指定周期模的多复变解析函数与第二库桑问题（Analytische Funktionen mehrerer komplexer Veränderlichen zu vorgegebenen Periodizitätsmoduln und das zweite Cousinsche Problem）。Math. Ann. 123 (1951), 201--222。
\item[{[63]}] Thimm, W.：阿贝尔函数代数相关的魏尔斯特拉斯定理及其推广（Der Weierstraßsche Satz der algebraischen Abhängigkeit von Abelschen Funktionen und seine Verallgemeinerungen）。载 Festschrift zur Gedächtnisfeier für Karl Weierstrass 1815--1965, H. Behnke、K. Kopfermann 编，Köln--Opladen: Westdeutscher Verlag, 1966。
\end{description}
